\chapter{Metodología}
\label{cap:metodologia}

\section{Herramientas de Simulación y Procesado}
\label{sec:metodologia:herramientas}

\subsection{Framework de simulación: penRed.}
% qué es penRed
% Descripción del software y su relación con PENELOPE (motor físico).
% Características principales: Modularidad, paralelismo y lenguaje C++.
% El módulo PENNUC: Justificación de su uso para la simulación de decaimientos complejos (\yochoseis)
\subsubsection{Definiciones Clave de la Simulación}
% Concepto de Historia (History): Definición como evento de decaimiento único y su importancia para la validación ("Same History").
% Concepto de Singles vs. Coincidencias: Diferencia entre el dato crudo (output de penRed) y el dato procesado.
% Tallies: Explicación del Tally Singles utilizado para la adquisición de datos en modo lista.

\subsection{Entorno de simulación: Cluster UPV}
\subsection{Software de reconstrucción: Bruker RecoServer (Algoritmo MLEM).}
\subsection{Análisis de imagen: AMIDE.}
% rois y obtención de la intensidad

\section{Definición del Escenario de Simulación}
\label{sec:metodologia:escenario_sim}

\subsection{Modelado del escáner: Geometría del Bruker PET/CT Si78.}
\subsection{El Fantoma Derenzo: Especificaciones y materiales.}
\subsection{Configuración de la fuente radiactiva (Archivos .nuc).}

\section{Desarrollo del Algoritmo de Post-Procesado (C++)}
\label{sec:metodologia:postprocess}

\subsection{Arquitectura del Software processs singles}
\subsection{Modelado de la respuesta del detector (Energy Blurring).}
\subsection{Implementación de algoritmos de coincidencia:}
\subsection{Same History (Validación ideal).}
\subsection{Closest Time (Simulación realista).}
\subsection{Multiplexing Mode (Detección de triples).}

\section{Diseño Experimental: Fases de Simulación}
\label{sec:metodologia:fases}

\subsection{Fase 1: Validación Ideal (Ground Truth).}
\subsection{Fase 2: Estudio de Efectos Físicos (Scatter y Atenuación).}
\subsection{Fase 3: Simulación Realista (Randoms).}
\subsection{Fase 4: Optimización de Ventanas (Energía y Tiempo).}
\subsection{Fase 5: Prueba de concepto de Multiplexado.}

\section{Métricas de Evaluación de Calidad}
\label{sec:metodologia:metricas}

\subsection{Definición de ROIs.}
\subsection{Coeficiente de Recuperación de Contraste (CRC).}
\subsection{Relación Señal-Ruido (SNR).}